% ============================================================================
% Fichier  : partie4/P04_C17_nuage.tex
% Titre    : Tester le Nuage et les Conteneurs
% Auteur   : MINKA MI NGUIDJOÏ Thierry Emmanuel
% Version  : 1.0 -- novembre 2026
% ============================================================================
% Sources : CIS Benchmarks (fournisseurs de nuage, Docker, Kubernetes) ;
% modèle de responsabilité partagée ; MITRE ATT&CK (matrices nuage et
% conteneurs). Services en nuage de LiZbiZ : cas fictif.
% ============================================================================

\chapter{Tester le Nuage et les Conteneurs}
\label{chap:nuage}
\chaptermark{Tester le nuage et les conteneurs}

\epigraph{%
    « Dans le nuage, la faille n'est presque jamais celle du fournisseur~;\\
    c'est celle de la configuration du client. »%
}{Principe du cours}

\niveauChapitre{Expert}{Chapitre~\ref{chap:api}~; notions
d'administration système.}

\begin{boxobjectifs}
\begin{itemize}
    \item Expliquer le modèle de responsabilité partagée.
    \item Reconnaître les erreurs de configuration les plus courantes dans
          le nuage.
    \item Comprendre le risque des secrets et des droits trop larges.
    \item Décrire les faiblesses des conteneurs et de leur orchestration.
    \item Tester un environnement en nuage dans les limites du contrat avec
          le fournisseur.
\end{itemize}
\end{boxobjectifs}

\bigskip

% ============================================================================
\section*{Question directrice : à qui appartient la faille quand le service est hébergé chez un fournisseur~?}
% ============================================================================

LiZbiZ héberge son nouveau portail et sa base chez un fournisseur de nuage.
M.~KOUNOU pense que la sécurité « est comprise dans l'abonnement ». Le test
montre le contraire~: le fournisseur sécurise son socle, mais la
configuration, les droits et les secrets restent la responsabilité de
LiZbiZ, et c'est là que sont les faiblesses.

\medskip
\textbf{Plan du chapitre.}
\begin{enumerate}
    \item La responsabilité partagée (section~\ref{sec:c17-partage}).
    \item Les erreurs de configuration (section~\ref{sec:c17-config}).
    \item Secrets et droits (section~\ref{sec:c17-droits}).
    \item Conteneurs et orchestration (section~\ref{sec:c17-conteneurs}).
    \item Tester dans les limites du contrat (section~\ref{sec:c17-limites}).
\end{enumerate}

\separateur

% ============================================================================
\section{La Responsabilité Partagée}
\label{sec:c17-partage}
% ============================================================================

\begin{boxdef}
\termecle{Modèle de responsabilité partagée}~: répartition des tâches de
sécurité entre le fournisseur de nuage (socle matériel, virtualisation,
services de base) et le client (configuration, droits, données,
applications). La frontière dépend du service souscrit.
\end{boxdef}

% ============================================================================
\section{Les Erreurs de Configuration}
\label{sec:c17-config}
% ============================================================================

\begin{tableaucours}{L{4.4cm}L{7.8cm}}{%
    \tableauHeader{Erreur fréquente} &
    \tableauHeader{Conséquence}
}
Espace de stockage ouvert au public & Fichiers et sauvegardes lisibles par
tous \\
Service d'administration exposé & Console ou base accessible depuis
l'Internet \\
Journalisation désactivée & Aucune trace en cas d'incident \\
Chiffrement non activé & Données au repos en clair \\
Réseau trop ouvert & Règles de filtrage permissives \\
\end{tableaucours}
\vspace{-0.8em}
\begin{center}{\small\itshape Tableau~17.1. Erreurs de configuration courantes}\end{center}

% ============================================================================
\section{Secrets et Droits}
\label{sec:c17-droits}
% ============================================================================

Deux risques dominent le nuage~: les \termecle{secrets} (clés d'accès, jetons)
laissés dans le code, les images ou les dépôts~; et les \termecle{droits}
trop larges accordés à un compte ou à une machine, qui permettent, depuis
une ressource compromise, d'atteindre tout l'environnement.

\begin{boxalert}
\textbf{Violation typique.} Utiliser une clé d'accès trouvée pour agir
au-delà du périmètre.

\smallskip
\textbf{Règle~:} une clé découverte se signale et se démontre dans la
limite stricte du périmètre convenu. On ne crée pas de ressources, on ne
touche pas aux données des autres clients du fournisseur, on respecte les
règles de test de celui-ci.
\end{boxalert}

% ============================================================================
\section{Conteneurs et Orchestration}
\label{sec:c17-conteneurs}
% ============================================================================

\begin{tableaucours}{L{4.4cm}L{7.8cm}}{%
    \tableauHeader{Faiblesse} &
    \tableauHeader{Risque}
}
Image non à jour & Vulnérabilités connues embarquées \\
Conteneur privilégié & Sortie du conteneur vers l'hôte \\
Secret dans l'image & Clé lisible par qui tire l'image \\
Tableau de bord d'orchestration exposé & Contrôle du groupe de conteneurs \\
Absence de limites & Un conteneur épuise l'hôte \\
\end{tableaucours}
\vspace{-0.8em}
\begin{center}{\small\itshape Tableau~17.2. Faiblesses des conteneurs}\end{center}

% ============================================================================
\section{Tester dans les Limites du Contrat}
\label{sec:c17-limites}
% ============================================================================

\begin{boxprocedure}
\textbf{Avant un test en nuage}
\begin{enumerate}
    \item Lire les règles de test du fournisseur~; certaines actions sont
          interdites ou exigent une déclaration.
    \item Délimiter les comptes et ressources du périmètre.
    \item Convenir des actions permises (lecture de configuration, oui~;
          création de ressources, non).
    \item Prévoir la conduite en cas de découverte d'un secret.
    \item Vérifier que l'environnement testé n'est pas partagé avec d'autres
          clients.
\end{enumerate}
\end{boxprocedure}

\begin{boxobjet}{T-10 · Service de test exposé sur l'Internet}
\textbf{Description~:} une interface de programmation de test, hébergée dans
le nuage, est accessible publiquement sans authentification.
\textbf{Preuve~:} adresse et réponse obtenue (données masquées).
\textbf{Gravité~:} moyenne. \textbf{Impact~:} divulgation de données de test,
parfois proches des données réelles~; point d'entrée. \textbf{Recommandation~:}
restreindre l'accès, authentifier, retirer les environnements de test
inutiles. \textbf{Détection~:} inventaire périodique des services exposés.
\end{boxobjet}

\begin{boxmetier}{Auditeur de sécurité et RSSI}
\textbf{Auditer la configuration, pas le fournisseur.} L'audit du nuage
porte sur ce que le client configure~: droits, secrets, exposition,
journalisation. Les référentiels de configuration sécurisée (CIS) donnent la
liste des points à vérifier.
\end{boxmetier}

\separateur

\begin{boxresume}
\textbf{Ce qu'il faut retenir de ce chapitre~:}
\begin{itemize}
    \item \textbf{La responsabilité est partagée}~; la faille est presque
          toujours côté client.
    \item \textbf{Les erreurs de configuration} (stockage ouvert, service
          exposé) dominent.
    \item \textbf{Secrets et droits trop larges} sont les deux grands
          risques.
    \item \textbf{Les conteneurs} ajoutent leurs faiblesses propres.
    \item \textbf{Le test en nuage} respecte les règles du fournisseur.
\end{itemize}
\end{boxresume}

\begin{boxinfo}
\textbf{Points de contrôle de l'élève.} Je sais~: expliquer la
responsabilité partagée~; citer les erreurs de configuration~; expliquer le
risque des secrets et des droits~; décrire les faiblesses des conteneurs~;
préparer un test en nuage conforme.
\end{boxinfo}

\bigskip

\begin{boxexo}[title={\ding{45}\quad Exercice~17.1. Frontière \niveaubase}]
Qui, du fournisseur ou du client, est responsable de la configuration des
droits d'accès~?
\end{boxexo}

\begin{boxexo}[title={\ding{45}\quad Exercice~17.2. Stockage ouvert \niveaubase}]
Pourquoi un espace de stockage ouvert au public est-il dangereux~?
\end{boxexo}

\begin{boxexo}[title={\ding{45}\quad Exercice~17.3. Secret \niveaubase}]
Citez trois endroits où un secret peut traîner dans un environnement en
nuage.
\end{boxexo}

\begin{boxexo}[title={\ding{45}\quad Exercice~17.4. Droits larges \niveaumoyen}]
Expliquez comment des droits trop larges transforment une ressource
compromise en compromission de tout l'environnement.
\end{boxexo}

\begin{boxexo}[title={\ding{45}\quad Exercice~17.5. Conteneur privilégié \niveaumoyen}]
Pourquoi un conteneur privilégié met-il en danger l'hôte~?
\end{boxexo}

\begin{boxexo}[title={\ding{45}\quad Exercice~17.6. Clé trouvée \niveaumoyen}]
Vous découvrez une clé d'accès valide. Que faites-vous, et que ne faites-vous
pas~?
\end{boxexo}

\begin{boxexo}[title={\ding{45}\quad Exercice~17.7. T-10 \niveaumoyen}]
Rédigez le plan de correction de T-10.
\end{boxexo}

\begin{boxexo}[title={\ding{45}\quad Exercice~17.8. Règles du fournisseur \niveaumoyen}]
Pourquoi faut-il lire les règles de test du fournisseur avant un test en
nuage~?
\end{boxexo}

\begin{boxexo}[title={\ding{45}\quad Exercice~17.9. Audit de configuration \niveauexpert}]
Construisez une liste de contrôle de dix points pour auditer la
configuration en nuage de LiZbiZ.
\end{boxexo}

\begin{boxexo}[title={\ding{45}\quad Exercice~17.10. Responsabilité \niveauexpert}]
Rédigez une note qui clarifie, service par service, la frontière de
responsabilité entre LiZbiZ et son fournisseur.
\end{boxexo}

\bigskip

\begin{boxinfo}
\textbf{Pour aller plus loin.}
\begin{itemize}
    \item CIS Benchmarks~\cite{cis2021controls}~; MITRE ATT\&CK~\cite{mitreattack}.
\end{itemize}
\end{boxinfo}

\begin{boxinfo}
\textbf{Transition.} La Partie~V ouvre sur les tests au-delà du système
d'information classique~: le Chapitre~\ref{chap:sansfil} traite du sans-fil
et des objets connectés.
\end{boxinfo}
